\chapter{Conclusão}\label{cap:conclusao}

Este trabalho avaliou a detecção de falhas por afastamento da operação normal com um sistema de baixo custo que adquire e transmite o sinal bruto de vibração, em \campoEnsaios\ ensaios de \campoMaquinas\ máquinas de uma instalação industrial e numa base pública com desbalanceamento real. A cadeia de aquisição entregou completas as \campoJanelas\ janelas dos trechos analisados, com o resumo SHA-256 conferido, e o valor eficaz da aceleração nas máquinas em operação ficou de \sinalRazaoMin\ a \sinalRazaoMax\ vezes o do sensor em repouso. A detecção não atendeu aos critérios registrados, e nenhuma das quatro hipóteses foi corroborada.

O sistema construído para a coleta tem um transmissor, um \textit{gateway}, uma API com banco de dados e um aplicativo de mesa, configurados sem cabo depois da primeira gravação. Nos ensaios de ponta a ponta, ele cumpriu a medição sem enlace, a retenção dos dados com a API fora do ar, o envio do sinal bruto e a atualização sem cabo do \textit{gateway}. Os \sisTestesCampoCasos\ casos de teste automatizados aplicados às funções dos \textit{firmwares} da coleta passaram. Na coleta, todas as janelas dos trechos analisados chegaram completas, e, fora do período de reinícios, chegaram completas \rajadasForaPct~\% das rajadas do dia. O envio do sinal bruto permitiu fazer seis análises posteriores sobre as mesmas janelas, o que não seria possível num nó que envia só a decisão do detector.

Sem medição da tensão da bateria, que não tinha sido carregada antes da coleta, o transmissor reiniciou a cada ciclo por queda de tensão, e a segunda rodada da dobradeira ficou sem sinal bruto. O consumo do transmissor também não foi medido. O retorno automático da configuração foi verificado só no computador, e o canal cifrado não cobre a configuração inicial do \textit{gateway}. Para ficar instalado por meses, o transmissor ainda precisa de atualização sem cabo, de um gravador para as janelas medidas sem o \textit{gateway} e de autenticação dos comandos recebidos pelo rádio.

O falso alarme entre rodadas (H1) ficou indeterminado, com \hUmCalandraIFDireto~\% das janelas da segunda rodada da calandra marcadas e intervalo de confiança de \hUmCalandraIFDiretoInf~\% a \hUmCalandraIFDiretoSup~\%. A menor mudança detectável em $d_1$ (H2) passou de \SI{0,7}{\milli\meter\per\second} nas duas rodadas da calandra, com \hDoisMmdCalandraUm\ e \hDoisMmdCalandraDois~mm/s, e no compressor, com \hDoisMmdCompressorUm~mm/s. Com a assinatura de desbalanceamento somada a janelas reais (H3), o detector principal não marcou metade das janelas até \SI{2,8}{\milli\meter\per\second} em nenhuma máquina, e a maior fração marcada foi de \hTresFracaoMaxCalandra~\%, na calandra. Na base pública (H4), ele marcou de \hQuatroFracaoTres~\% a \hQuatroFracaoQuatro~\% das janelas com desbalanceamento real e \hQuatroFracaoZero~\% das janelas sem desbalanceamento medidas depois da remontagem do suporte.

As análises posteriores mostraram quais desses resultados dependem do detector e da amostra. Com as sementes de 0 a \num{99} do gerador de números aleatórios, H1 seria corroborada com \apTresVereditoUmCorroborada\ das \apTresSementes, e a mediana do falso alarme da calandra foi de \apTresSementeCalandraMed~\%, enquanto H4 seria refutada com \apTresVereditoQuatroRefutada. A refutação de H2 se manteve em todas as réplicas da reamostragem da calandra e do compressor, e a menor mudança detectável dos cinco ensaios do critério ficou de \apCincoRazaoMin\ a \apCincoRazaoMax\ vezes a do sensor em repouso. A fase sorteada da assinatura não explica a refutação de H3. Com os descritores que a assinatura altera além da faixa do treino, o escore do Isolation Forest ficou abaixo do limiar na guilhotina, na calandra e no compressor. A mediana de até cinco janelas por medição não levou a menor mudança detectável abaixo de \SI{0,7}{\milli\meter\per\second} na calandra e no compressor.

Para quem instala um nó de baixo custo, os limites observados estão no detector, na linha de base e na montagem do sensor. O ruído do sensor responde por uma parte pequena da menor mudança detectável, cerca de \SI{5}{\percent} no ensaio com a menor razão, se as variações do repouso e das máquinas forem independentes (seção~\ref{sec:disc-sensibilidade}). O Isolation Forest marca uma janela que se afasta do treino em vários descritores ao mesmo tempo, e o escore dele não cresce além da faixa do treino. A regra por mediana e MAD, que avalia cada descritor separadamente, marcou metade das janelas com a assinatura em fase com \apUmMadAcinquentaZeroDobradeira~mm/s na dobradeira e \apUmMadAcinquentaZeroGuilhotina~mm/s na guilhotina, valores ainda acima do critério de \SI{0,7}{\milli\meter\per\second}.

Com a linha de base de uma rodada, de \campoJanelasGuilhotinaUm\ a \campoJanelasCalandraUm\ janelas por máquina, o falso alarme estimado variou com a semente do detector a ponto de mudar a classificação de H1. O detector marcou operação normal medida depois de mudanças simultâneas na fixação, na alimentação, no intervalo entre ciclos e na temperatura do sensor, e os dados não indicam qual dessas mudanças produziu a marcação. A variação entre janelas medida nas máquinas reúne a vibração da máquina e a resposta da fixação do transmissor, que não foi medida.

Entre os trabalhos consultados, em busca não exaustiva, este trabalho acrescenta quatro resultados sobre a detecção com nós de baixo custo. O primeiro é a avaliação em campo, entre rodadas da mesma máquina, de um detector treinado só com operação normal, com hipóteses registradas antes do cálculo e resultado negativo relatado. O segundo é a medição do teto do escore do Isolation Forest fora da faixa do treino. O terceiro é a medida da variação do falso alarme com a semente do detector e com o treino. O quarto é um sistema de aquisição de baixo custo que transmite o sinal bruto, verificado diante de requisitos definidos antes da coleta.

Os resultados levam a três predições que uma coleta nova nas mesmas máquinas pode testar. A primeira é que, com o treino formado por janelas de mais de uma rodada e de mais de uma fixação, a mediana do falso alarme entre as sementes de 0 a \num{99}, na rodada seguinte da calandra, fica abaixo dos \apTresSementeCalandraMed~\% obtidos com o treino de uma rodada. A segunda é que, se a resposta da fixação explica as medianas axiais acima do limite entre as zonas C e D, um sensor fixado de forma rígida no mesmo ponto mede valores menores na dobradeira, na calandra e no compressor. A terceira é que, no compressor, a menor mudança detectável só diminui com medições que cubram vários períodos da variação lenta indicada pela autocorrelação de \apDoisAutocorrCompressorUm\ entre janelas separadas por \campoCicloUm~s, período que a coleta não mediu.

Medir a detecção de falhas reais em campo exige máquinas com falha conhecida, ou falhas semeadas em máquinas da instalação, acompanhadas por semanas, com um sensor de referência no mesmo ponto e o registro da carga e da temperatura. Para defeitos de rolamento, exige também um sensor com banda acima dos \SI{184}{\hertz} nominais do filtro usado, porque a análise de envelope trabalha na faixa das ressonâncias estruturais (seção~\ref{sec:disc-banda}).
